\documentclass[10pt]{article}

\usepackage[utf8]{inputenc}

\usepackage[T1]{fontenc}

\usepackage{amsmath}

\usepackage{amsfonts}

\usepackage{amssymb}

\usepackage[version=4]{mhchem}

\usepackage{stmaryrd}

\usepackage{graphicx}

\usepackage[export]{adjustbox}

\graphicspath{ {./images/} }



\begin{document}

бесконечно малого изгибания с р е д и в в о й п ов е р х в о с т и, соответствующей $F$ и $F^{*}$ (см. КонФоссена преобразование) (при условии $|\chi|<\pi$ ова регулярна).



Свойства П. Д. для ивометричных поверхностей положительвой гауссовой кривизны аналогичны свойствам индикатрисы в ращений: напр., удельная внутренняя кршвизна П. Д. всюду отрицатөльна, и потому она иғрает при исследовании изометрии выпуклых поверхностей такую же роль, что и индикатриса вращевий. М. И. Войцеховский.



ПОВТОРНОГО ЛОГАРИФМА ЗАКОН - предельная теорема теории вероятностей, являющаяся уточнением больших чисел усиленного закона. Пусть $X_{1}, X_{2}, \ldots-$ последовательғость случайных величин и



\[

S_{n}=X_{1}+X_{2}+\ldots+X_{n}

\]



Для простоты предполагается, что $S_{n}$ для каждого $n$ имеет нуль своей медианой. В то время как теоремы об усиленном законе больших чисел указывают условия, при к-рых $\frac{s_{n}}{a_{n}} \rightarrow 0$ почти наверное (п. н.) при $n \rightarrow \infty$, где $\left\{a_{n}\right\}-$ числовая последовательность, теоремы о П. л. З. имеют дело с числовыми последовательностями $\left\{c_{n}\right\}$ такими, что



\[

\limsup _{n \rightarrow \infty} \frac{s_{n}}{c_{n}}=1 \text { п.н. }

\]



или



\[

\limsup _{n \rightarrow \infty} \frac{\left|S_{n}\right|}{c_{n}}=1 \text { п.П. }

\]



Соотношение (1) равносильно тому, что



$\boldsymbol{\Pi}$



\[

P\left\{S_{n}>(1+\varepsilon) c_{n} \text { б.ч. }\right\}=0

\]



\[

\mathrm{P}\left\{S_{n}>(1-\varepsilon) c_{n} \text { б.ч. }\right\}=1

\]



для любого $\varepsilon>0$, где запись «б. ч.» означает бесконечное число раз.



Соотношения вида (1) и (2) справедливы при более ограничительных условиях, чем оценки, вытекающие из усиленного закона больших чисел. Если $\left\{X_{n}\right\}-$ последовательность независимых случайных величин, имеющих одинаковое распределение с математич. ожиданием, равным нулю, то



\[

\frac{s_{n}}{n} \rightarrow 0 \text { п.н. при } n \rightarrow \infty

\]



(т өо рем а К о мо го ро в а); если выполнено дополнительное условие $0<\mathrm{E} X_{1}^{2}<\infty$, то имеет место более сильное соотношение (2), в к-ром



\[

c_{n}=(2 n b \ln \ln (n b))^{1 / 2}

\]



где $b=E X_{1}^{2}($ те о рем а X а р т а а а - Вин тн е p a).



Первой теоремой общего типа о законе повторного логарифма был следующий результат А. Н. Колмогорова [1]. Пусть $\left\{X_{n}\right\}$ - последовательность независимых случайных величин с математич. ожиданиями, равными нулю, конечными дисперсиями и пусть



\[

B_{n}=\sum_{k=1}^{n} E X_{k}^{2}

\]



Если $B_{n} \rightarrow \infty$ при $n \rightarrow \infty$ п существует последовательность положительных постоянных $\left\{M_{n}\right\}$ такая, что



\[

\left|X_{n}\right| \leqslant M_{n} \text { и } M_{n}=o\left(\left(\frac{B_{n}}{\ln \ln B_{n}}\right)^{1 / 2}\right)

\]



то выполнены соотношения (1) и (2) при



\[

c_{n}=\left(2 B_{n} \ln \ln B_{n}\right)^{1 / 2} .

\]



В частном случае, когда $\left\{X_{n}\right\}-$ последовательвость независимых случайных величин, имеющих одинако- вое распределение с двумя значениями, это утверждение было получено А. Я. Хинчиным [2]. Ю. Марцинкевич и А. Зигмунд [3] показали, что в условиях теоремы Колмогорова нельзя заменить $о$ на $O$. Обобщения П. л. з. Колмогорова для последовательностей независимых ограниченных неодинаково распределенных случайных величин были исследованы В. Феллером [4]. Другие обобщения П. л. з. см. в [5]; имеется также следующий результат (см. [6]), примыкающий к теорөме Хартмана - Винтнера: если $\left\{X_{n}\right\}$ - последовательность независимых случайных величин, имеющих одинаковое распределение с бесконечной дисперсией, то



\[

\limsup _{n \rightarrow \infty} \frac{\left|S_{n}\right|}{(n \ln n \ln n)^{1 / 2}}=\infty \text { П.н. }

\]



Результаты, полученные в области П. л. 3. для последовательностөй независимых случайных величин, послужили отправным пунктом для многочисленных исследований применимости П. л. З. к последовательностям зависимых случайных величин и векторов и к случайным процессам.



Лum.: [1] K о Л M о г о p о в A. H., "Math. Ann.", 1929, Bd 101, S. 126-35; [2] X и н ч и н A. A., "Fundam. math.", 1924,

v. 6, p. 9-20; [3] M a r c i n k i e w i c z J., Z y g un d A., там же, 1937, v. 29, p. 215-22; [4] F e 11 e r W., "Trans. Amer. Math. Soc.", 1943, v. 54, p. 373-402; [5] St t a s e en V., "Z. Wahr[6] е г о ж е, там же, 1965-66, Bd 4, S. 265-68; [7] H a r t-



\includegraphics[max width=\textwidth, center]{2023_07_08_d0c95becefb217498745g-1}

[9]'П е т р о в В. В., Суммы независимых случайных величин, М., 1972.



ПОВТОРНЫЙ ИНТЕГРАЛ - интеграл, в к-ром последовательно выполняется интегрирование по разным переменным, т. е. интеграл вида



\[

\int_{A_{y}}\left[\int_{A_{(y)}} f(x, y) d x\right] d y .

\]



Функция $f(x, y)$ определена на множестве $A$, лежащем в прямом произведении $X \times Y$ пространств $X$ и $Y$, в к-рых заданы $\sigma$-конечные меры $\mu_{x}$ и $\mu_{y}$, обладающие свойством полноты; множество $A(y)=\{x:(x, y) \in A\} \subset X$ («сечение» множества $A$ ), измеримое относительно меры $\mu_{x} ;$ множество $A_{y}$ (проёцця множества $A$ в пространство $Y$ ), измеримое относительно меры $\mu_{y}$. Интегрирование по $\boldsymbol{A}(y)$ производится по мере $\mu_{x}$, а по $A_{y}$ - по мере $\mu_{y}$. Интеграл (1) обозначают также



\[

\int_{A_{\boldsymbol{y}}} d y \int_{\boldsymbol{A}(y)} f(x, y) d x .

\]



К П. и. могут быть сведены кратные интегралы.



Пусть функция $f(x, y)$, интегрируемая по мере $\mu=\mu_{x} \times \mu_{y}$ на множестве $A \subset X \times Y$, продолжена нулем на все пространство $X \times Y$, тогда П. и.



\[

\int_{\boldsymbol{Y}} d y \int_{\boldsymbol{X}} f(x, y) d x

\]



И



\[

\int_{\boldsymbol{X}} d x \int_{\boldsymbol{Y} !} f(x, y) d y

\]



существуют и равны между собой:



\[

\int_{\boldsymbol{Y}} d y \int_{\boldsymbol{X}} f(x, y) d x=\int_{\boldsymbol{X}} d x \int_{\boldsymbol{Y}} f(x, y) d y

\]



(см. Фубини теорема). В левом интеграле внешнее интегрирование фактически производится по множеству $A_{y}^{*}=\left\{y: y \in A_{y}, \mu_{x} A(y)>0\right\}$. Таким образом, в частности, для точек $y \in A_{y}^{*}$ множества $A(y)$ измеримы относительно меры $\mu_{x}$. По всему множеству $A_{y}$ брать этот интеграл, вообще говоря, нельзя, т. к. при измеримом относительно меры $\mu$ множества $A$ множество $A_{y}$ может оказаться неизмеримым относительно меры $\mu_{y}$, так же, как и отдельные множества $\boldsymbol{A}(y), y \in A_{y}$,





\end{document}